% Ablations: what a modification is, the table of paired comparisons and the
% two figures of validation curves. The same for every task, so a model paper
% inputs this file at the start of its ablation section and then reports its
% own findings. Each part appears once the task has the runs it needs, and the
% list below needs at least one of its items.
% The prose follows docs/writing.md.
\ifpairs
The main results compare complete models. An ablation isolates one factor: it
compares a model with its \emph{base}, a model that differs from it in a
single respect. Table~\ref{tab:pairs} lists every such comparison for this
task, with the gain of the modified model over its base on the validation set
and its paired gain on the test set, computed on the same resampled test
games as in the main results. The acceptance criterion of
Section~\ref{sec:protocol} is applied to these pairs. The modifications are
the following.
\begin{itemize}
\ifchangehand
\item \textbf{Hand-crafted features} (suffix \code{-hand}): the network
receives the hand-crafted features as additional context features.
\fi
\ifchangeepochs
\item \textbf{Longer schedule} (suffix \code{-long}): twice the training
epochs, with the same warm-up fraction and cosine schedule.
\fi
\ifchangestart
\item \textbf{Pretraining} (suffix \code{-pretrained}): the transformer
encoder is initialised from a transformer of the same size that was trained
to predict the engine's win probability for positions from the Lichess
evaluation database~\cite{lichessdb}. The piece and square embeddings and all
transformer layers are copied. The prediction head is initialised randomly,
as are the projections of the auxiliary planes that the pretraining task does
not have and, unless the two tasks have the same context features, the
projection of the context features. Training then follows the standard
schedule.
\fi
\ifchangeattacks
\item \textbf{Attack-aware attention bias} (suffix \code{-attacks}): the
relative position bias encodes where two squares lie and not what stands
between them. With this modification each attention head learns four
additional biases for a pair of squares, applied when the piece on one square
attacks the other, when it is attacked from the other, and for the same two
relations through exactly one intervening piece, as in a pin or a skewer. The
relations are computed from the tokens without search, and the biases are
initialised to zero.
\fi
\ifchangegeometry
\item \textbf{No relative position bias} (suffix \code{-no-geometry}): the
relative position bias is removed. This configuration measures the
contribution of the bias and is not a candidate model.
\fi
\ifchangeconv
\item \textbf{ResNet encoder} (prefix \code{conv-medium}): a ResNet with
squeeze-and-excitation in place of the transformer, trained at the higher
peak learning rate used for the convolutional models.
\fi
\ifchangehybrid
\item \textbf{Hybrid encoder} (prefix \code{hybrid-medium}): a few residual
convolutional blocks followed by transformer layers over the resulting
square representations.
\fi
\ifchangelarger
\item \textbf{Larger transformer} (prefix \code{transformer-large}): more layers and
a larger width, with the training schedule unchanged.
\fi
\ifgames
\item \textbf{Training-set size}: the same model trained on a fraction or a
multiple of the training games (Section~\ref{sec:training-size}).
\fi
\iffeatures
\item \textbf{Feature groups}: a model trained without one group of
hand-crafted features at a time (Section~\ref{sec:feature-groups}).
\fi
\end{itemize}

\longblock{pairs}{Paired comparisons. Each row compares a model with a
baseline, in most rows the model it modifies. \emph{Change} states how the
two differ. The gain is that of the model over the baseline and is positive
when the model is better. On
the validation set it is the difference between the two scores. On the test
set it is the paired gain, with its 95\% CI, and is set in bold when
statistically significant.}{tab:pairs}
\fi

\ifvariants
Figure~\ref{fig:variants} shows the validation curves of each model of the
fixed sequence whose training schedule or initialisation was modified,
together with the modified models. A modified model has the colour of its
base and a broken line: long dashes for a longer schedule and dots for a
pretrained encoder.

\begin{figure}[htbp]\centering
\block{variants}
\caption{Validation curves of a model of the fixed sequence (solid line) and
of the models that modify its training schedule or initialisation (broken
lines), after every epoch. One panel per base model, on a common scale. The
dashed horizontal line is the GBDT baseline's validation score.}
\label{fig:variants}
\end{figure}
\fi

\ifshapes
Figure~\ref{fig:shapes} shows the validation curves of the models that
change the architecture and nothing about the training data, each in its own
panel beside its base.

\begin{figure}[htbp]\centering
\block{shapes}
\caption{Validation curves of a model with a modified architecture (broken
line) and of its base (solid line), after every epoch. One panel per modified
model, on a common scale. The dashed horizontal line is the GBDT baseline's
validation score.}
\label{fig:shapes}
\end{figure}
\fi
